\documentclass[11pt,a4paper]{article}
\usepackage[T1]{fontenc}
\usepackage[utf8]{inputenc}
\usepackage{lmodern,microtype}
\usepackage[margin=0.85in]{geometry}
\usepackage{amsmath,amssymb,amsthm,booktabs,tabularx,array,enumitem}
\usepackage[hidelinks]{hyperref}
\setlength{\parindent}{0pt}
\setlength{\parskip}{5pt}
\setlength{\emergencystretch}{2em}
\setlist{nosep,leftmargin=*}
\newcolumntype{Y}{>{\raggedright\arraybackslash}X}
\newcommand{\Tag}[1]{\ensuremath{\mathsf{#1}}}
\newcommand{\Proc}[1]{\mathsf{#1}}
\newcommand{\Hash}{\mathsf H}
\newtheorem{proposition}{Proposition}
\newtheorem{assumption}{Assumption}
\hypersetup{pdftitle={Staking and Rotation Module v1}}
\begin{document}
\begin{center}
{\Large\bfseries Staking and Rotation Module}\\[4pt]
{\large Version 1: initial modular design}\\[3pt]
10 October 2026
\end{center}

\textbf{Purpose.} This module manages stake-backed membership, derives
regional and local leader orders, and applies multisignature regional
eligibility updates. A regional BFT engine accepts a fixed weighted member
set and a function $v\mapsto\ell(v)$, and implements its own view changes.
We do not change its voting, commitment, or timeout rules. The initial
profile uses one collateral token, fixed region identifiers, whole-position
joins and exits, and one occurrence of each identity per local cycle.
Partial exits, delegation, region migration, and stake-proportional leader
frequency are extensions, not implicit features.

\textbf{Design status.} This is the first candidate in a design--review--patch
sequence. The construction follows the fixed-snapshot and status/closure
separation in \texttt{main\_v3.tex}. It deliberately serializes every
cross-region handoff on L1, including ordinary quota expiry. No L1
transaction is required for an ordinary batch or an in-region view change.

\section{State and external contracts}
\label{sec:interfaces}

L1 is an authenticated, finalized, totally ordered control log. Clients
accept control actions only after finality. Its block timestamp is the
control clock $t$; a local wall-clock observation changes no authority.
Actions in the same block follow transaction order. Due scheduled events
are processed before an action at their effective timestamp.

Let $\mathcal R$ be a fixed, canonically ordered set of regions. A stake
position $o$ has a unique identifier, an immutable region $r(o)$, public
key, withdrawal address, and initially deposited amount $b_o$. Keys cannot
be shared by two concurrent positions. Position IDs remain available for
historical verification. A configuration epoch $e$ contains a frozen table
$S_e=\{(o,r(o),pk_o,s_o^e)\}$ with positive weights and
\[
 V_r^e=\{o\in S_e:r(o)=r\},\qquad
 W_r^e=\sum_{o\in V_r^e}s_o^e.
\]
Thus the requested notation $s_{r,j}^{(e)}$ denotes $s_o^e$ for the
$j$th position in region $r$. Weight is the recorded stake at the snapshot,
not a continuously recomputed quorum denominator.

The registry stores $(e,S_e,t_e,\rho_e,n,\alpha,A,\chi,P)$: configuration,
its actual activation timestamp, seed, a monotone authority counter,
the current session or $\bot$, the last closed checkpoint, the effective
status map, and at most one frozen next-epoch plan. A session is
$\alpha=(e,n,r,A_{\rm start})$. The application height is global;
the local view starts at zero for every fresh session. All signatures bind
the chain, registry, configuration, and session, in addition to their
message-specific fields.

\begin{assumption}[Required service contracts]
\label{ass:services}
The regional BFT provides prefix safety and its stated conditional progress
for each fixed configuration and supplied cyclic leader map. Its fault bound
must hold separately in every serving region, measured using $S_e$.
The L1 control log is safe. Signature and hash checks are sound.
The interfaces below have their stated semantics.
\end{assumption}

\begin{tabularx}{\linewidth}{@{}p{.26\linewidth}Y@{}}
\toprule
Interface & Required behavior\\\midrule
$\Proc{Open}(\alpha,S_e,\ell,A)$ &
Start only from the finalized L1 authority for $\alpha$, after fetching
and verifying $A$ and its required data. Repeated calls restore the same
persisted session, never reset its signing state.\\
$\Proc{Final}(\alpha)$ &
Export a finalized application prefix with a publicly verifiable native
BFT finality witness under the historical configuration.\\
$\Proc{Close}(\alpha,\gamma)$ &
Order a terminal control command for the justified trigger $\gamma$;
return $(A^*,C)$, a checkpoint and terminal finality certificate.\\
$\Proc{VerifyClose}(\alpha,\gamma,A^*,C)$ &
Verify that $A^*$ contains every finalized or already-finalizable application
prefix of $\alpha$ and that no valid application continuation beyond
$A^*$ can subsequently arise in that session.\\
$\Proc{Seed}(e,d)$ &
Publish one verifiable seed for epoch $e$ and frozen-plan digest $d$.
It is unpredictable before the plan is fixed. Publication progress is
an additional liveness assumption.\\
\bottomrule
\end{tabularx}

\textbf{Closure is an additional BFT integration requirement.}
A terminal command must be ordered in the same log as application batches,
with application validity forbidding further application entries after it.
An ordinary checkpoint QC, signatures on the latest observed tip, or a
local stop timer does not implement this contract. If a failed region
cannot produce $C$, this minimal module waits. A more available recovery
mechanism must separately prove the same closure contract.

\section{Escrow, requests, and bounded membership changes}
\label{sec:stake}

For each region fix positive inbound/outbound stake budgets $Q_r^+$ and
$Q_r^-$, request-count budgets $N_r^+$ and $N_r^-$, and a minimum bond
$s_{\min}>0$. Each admitted position has
\[
s_{\min}\le b_o\le s_{\max,r}\le\min\{Q_r^+,Q_r^-\}.
\]
Fixed regional budgets also bound total changes by $\sum_r Q_r^\pm$.
Count gross additions and removals separately, even if net stake is zero.
The upper bound on an indivisible position prevents a permanent oversized
queue head. There is no permanent administrator-assigned membership seat.

\begin{enumerate}[label=\textbf{S\arabic*.},itemsep=4pt]
\item \textbf{Join.} Anyone deposits an admissible bond and registers the
position and proof of key possession. Its L1 sequence number determines
its order in region $r$'s inbound FIFO. It is \Tag{Pending}; it has no vote
or leader eligibility. After a fixed delay $T_J$ it can enter a snapshot.
The operator must synchronize before its first service session.
\item \textbf{Cancel pending.} The withdrawal owner may refund a pending
position only before it is selected into a frozen plan. Selected positions
are \Tag{Reserved}, cannot be refunded, and remain locked through activation.
\item \textbf{Request exit.} The owner of an active position may append one
irrevocable request to its region's outbound FIFO. A request does not stop
participation or release funds. The position remains active until actually
omitted by an installed successor configuration.
\item \textbf{Select changes.} At the cutoff, scan each FIFO in order.
Select the longest eligible prefix fitting both its stake and count budget;
leave the rest queued. Join eligibility includes the $T_J$ delay; exits
require membership in the current configuration. Different regions and
the two directions have independent queues. Apply exits and then joins to
form the next table. No minimum regional stake may veto an otherwise
eligible exit; an empty region simply becomes unselectable.
\item \textbf{Retire and release.} On actual installation of the successor
configuration, removed positions enter \Tag{Cooling} at time $t_o^{\rm ret}$.
They no longer vote in new sessions, but retain required old data for
$T_D$ and remain liable through $t_o^{\rm ret}+T_D+T_E^{\rm evid}$.
After that deadline, release only the residual bond not subject to a
timely unresolved penalty claim. Anyone may submit evidence through the
external adjudication interface; only the withdrawal address receives funds.
\end{enumerate}

\textbf{Historical liability.} Slashing adjudication is external to this
module. A debit reduces the remaining escrow, never a historical quorum
denominator. Each snapshot records the balance available at its cutoff;
subsequent debits do not redraw that snapshot or its schedule. The same
bond is not independently recoverable once per historical epoch: cumulative
payments cannot exceed its remaining funds. We claim preservation of
liability against voluntary withdrawal, not an invariant lower bound on
collateral after penalties. Timely disputes may delay release until resolved.

If L1 or safe retirement stalls indefinitely, active and reserved funds may
remain locked indefinitely. Unconditional bounded withdrawal would require
an additional, proved expiration/fencing mechanism. Churn limits do not
prove that the next region still meets the BFT fault bound.

\section{One next-epoch plan}
\label{sec:epoch}

Fix a target epoch duration $T_{\rm ep}$ (for example 30 minutes), cutoff
offset $0<T_F<T_{\rm ep}$, join delay $T_J$, and a randomness source whose
round is fixed in advance by the protocol. Actual configuration $e$ begins
at its finalized L1 installation timestamp $t_e$. Define
\[
 F_e=t_e+T_F,\qquad D_e=t_e+T_{\rm ep}.
\]
The cutoff uses the complete canonical L1 request prefix with timestamps
strictly smaller than $F_e$. The first finalized block crossing $F_e$
makes that prefix determinate; any caller can materialize its result.
Before processing a cancellation or other membership action at or after
$F_e$, the registry first materializes the plan. Late callers cannot move
the cutoff, reorder requests, or change the selected set.

Using S4, form
\[
P_{e+1}=(e+1,\Hash(S_e),F_e,S_{e+1},J_{e+1},X_{e+1}),\quad
d_{e+1}=\Hash(P_{e+1}).
\]
Reserve selected joins and mark selected departures; neither change is
yet effective. Obtain $\rho_{e+1}=\Proc{Seed}(e+1,d_{e+1})$ only after
the cutoff. One seed reference is accepted; no retry, cancellation, or
alternative seed is allowed because a participant dislikes the result.
The plan is ready when its table and seed are publicly verifiable.

At $D_e$, request closure of the then-authorized session. Installing $e+1$
requires both a ready plan and safe closure, or an already idle registry.
Until installation, $S_e$ remains the verification configuration.
Nominal expiry initiates retirement; it is not an automatic switch of
weights or permission for an alternative region to produce. Only one
next plan exists, so a delay does not create competing future configurations.
The next deadline is measured from the actual installation of $e+1$.

\section{Derived rotation maps and multisignature control}
\label{sec:maps}

Let $\Proc{Perm}$ be a fully specified deterministic shuffle using canonical
IDs, domain-separated randomness and unbiased integer sampling. Define
\[
\begin{aligned}
\mathcal L_R^e&=\Proc{Perm}(\mathcal R,
 \Hash(\texttt{REGION},e,\rho_e)),\\
\mathcal L_r^e&=\Proc{Perm}(V_r^e,
 \Hash(\texttt{LEADER},e,r,\rho_e)),\\
\ell_{e,r}(v)&=\mathcal L_r^e[v\bmod |V_r^e|],\qquad V_r^e\ne\varnothing.
\end{aligned}
\]
The BFT owns $v$ and decides when to advance it. The module never restarts
a view because an operator requests exit. This profile randomizes order;
it does not make identity frequency proportional to stake. An adversary
may split stake across identities subject to the bond limits; this does
not change its total vote weight, but can increase its leader frequency.

A region is eligible at control time $t$ iff
\[
\Proc{Eligible}(e,r,t)\iff W_r^e>0\ \land\ \chi_t(r)=\Tag{live}.
\]
At epoch installation or from idle, select the first eligible region in
$\mathcal L_R^e$. At a within-epoch handoff from $r$, scan the next positions
of that cyclic list once, excluding $r$. If no other region is eligible,
renew $r$ in a fresh session if $r$ is itself eligible; otherwise remain
idle at the closed checkpoint. This explicit same-region renewal is a
small extension of the waiting rule in v3.

The status committee has fixed keys and threshold $q_C$. It submits
\[
U_j=(\text{chain},\text{registry},j,\Hash(U_{j-1}),a_j,\chi_j)
\]
with $q_C$ distinct signatures. The registry checks consecutive versions,
the predecessor, strictly increasing effective timestamps, and
$a_j\ge t_{\rm include}+T_N$ for notice $T_N>0$.
Only the finalized installed update is authoritative. It sets future
eligibility, not stake weights, random seeds, application checkpoints, or
an arbitrary successor. Committee key changes are outside this profile.

When an effective update makes the current region off, latch a closure
request for that exact session. A later live update does not cancel it.
Updates to inactive regions only affect subsequent selection. A recovering
region must fetch the current checkpoint before producing. The label
\Tag{live} authorizes selection; it is not a proof of actual responsiveness
or of the regional BFT fault assumption. A corrupt committee can filter
the published random order or halt service, even though it cannot forge
a closure certificate.

\section{The single authority transition}
\label{sec:transition}

Each regional session admits at most $K$ fresh application batches above
its start anchor. An empty batch may be used if the application permits
idle progress. Local view changes, re-certifications and terminal commands
do not consume this quota. The BFT adapter enforces the application count.

\begin{enumerate}[label=\textbf{R\arabic*.},itemsep=4pt]
\item \textbf{Trigger.} Request $\Proc{Close}(\alpha,\gamma)$ when its
quota is exhausted, its region has a latched off event, or the configuration
deadline has arrived. $\gamma$ authenticates the applicable trigger.
Closure uses the existing BFT; it does not discard in-flight commitments.
\item \textbf{Seal.} Any relayer submits
$\Proc{Advance}(n,\alpha,\gamma,A^*,C)$ to L1. Require the expected counter
and session to equal registry state and verify the native closure.
When already idle, an advance references the stored anchor and needs no
new closure. Apply every due status update first.
\item \textbf{Choose configuration.} If $t\ge D_e$ and the next plan is
ready, install it and execute S5 for departing positions. If the deadline
has arrived but the plan is not ready, store the closed anchor and stay
idle. Otherwise keep $e$.
\item \textbf{Choose region.} After an epoch change use the first eligible
entry of the new list. Otherwise use the within-epoch successor rule.
The rule reads the status map effective at the L1 transition timestamp;
the relayer cannot supply a different target or observation time.
\item \textbf{Commit authority.} Atomically store $A^*$, consume the old
session and increment $n$. Record either a fresh
$\alpha'=(e,n,r',A^*)$ or $\bot$. Only after this transaction is finalized
may $\Proc{Open}$ start the successor. A later status update is evaluated
against that successor, never against an already consumed predecessor.
\end{enumerate}

Initialization supplies one checkpoint and the first authenticated stake
table and seed, with no predecessor session. Every later session consumes
exactly one closed predecessor or resumes from the sole stored idle anchor.
An accepted old receipt is checked with its original configuration and
remains part of the inherited prefix; it is not reweighted under the new
table. No consumer treats a locally proposed successor as authorized.

\section{Timeouts and unavailable dependencies}
\label{sec:timeouts}

\begin{tabularx}{\linewidth}{@{}p{.25\linewidth}Y Y@{}}
\toprule
Clock or event & Action & What it does not authorize\\\midrule
Native BFT timeout & Advance local view under its own protocol and
$\ell_{e,r}$. & Change region, members, or stake denominator.\\
Observed service silence $T_{\rm obs}$ & Submit an observation to the
committee; the committee may schedule off. & An automatic slash or takeover.\\
Status notice $T_N$ & Make the installed map effective in L1 order;
latch closure if needed. & Claim that wall-clock preparation or data recovery
has completed.\\
Epoch deadline $D_e$ & Request closure and attempt the next configuration.
& Activate without closure and a verified plan.\\
Closure retry timeout & Retry or alert; invoke native BFT recovery.
& Substitute a locally latest checkpoint for $C$.\\
Randomness or L1 delay & Wait; retain historical state and escrow.
& Redraw the seed or refund an active/reserved position.\\
Cooling deadline & Release residual unencumbered funds under S5.
& Erase timely pending claims or old verification records.\\\bottomrule
\end{tabularx}

Progress bounds require bounds on L1 inclusion/finality and dissemination,
seed publication, native BFT closure, checkpoint recovery and adjudication.
Without these bounds, the table specifies safe waiting, not a finite
failover or withdrawal deadline. A timestamp notice interval alone does
not imply that all nodes had that much usable preparation time.

\section{Security claims and their limits}

\begin{proposition}[Unique authority and prefix preservation]
Under Assumption~\ref{ass:services}, accepted sessions have a unique L1
order and their finalized application outputs form compatible prefixes,
including across epoch and regional changes.
\end{proposition}
\begin{proof}
The expected-counter check permits only one transition to consume the
current authority in the finalized L1 history. Selection is a function
of that transition's state. Inside a session, regional BFT safety gives
compatible prefixes. A valid close preserves every finalizable old prefix
and forbids a further old-session continuation. A successor starts only
from that anchor and finalized authority. Induction over transitions,
with idle steps preserving the anchor, proves the assertion.
\end{proof}

\begin{proposition}[Withdrawal does not erase timely liability]
Selected joins and active positions cannot voluntarily withdraw. A removed
position remains locked through its actual retirement, data-retention and
evidence intervals, and until timely pending claims are resolved.
\end{proposition}
\begin{proof}
S2 excludes reserved positions from cancellation. S3 leaves an exit request
active. S5 changes the state only upon actual configuration installation
and is the only release rule. Its time and pending-claim checks give the
stated property. This does not increase the remaining balance after slashes.
\end{proof}

\textbf{Conditional progress.} With eventual finalized control actions,
seed publication, data recovery, and native BFT closure, the module can
install a ready successor configuration and open an eligible region.
FIFO progress additionally needs fixed positive budgets and finitely many
earlier eligible requests. Claim completion needs eventual dispute
resolution. Neither randomized lists nor churn caps imply the per-region
fault assumption. Fair service allocation further depends on committee
eligibility decisions and operator behavior.

\textbf{Deliberate availability limit.} A multisignature off decision alone
cannot recover an unavailable old region. Obtaining a safely closed cut
without that region's cooperation remains an external recovery requirement.
This module formalizes that boundary rather than treating elapsed time
as a proof of closure.
\end{document}
